% Part: many-valued-logic
% Chapter: syntax-and-semantics
% Section: introduction

\documentclass[../../../include/open-logic-section]{subfiles}

\begin{document}

\olfileid{mvl}{syn}{int}

\olsection{جان پہچان}

کلاسیکی منطق وچ اسیں !!{formula}s نال کم کردے آں جیہڑے
!!{propositional variable}s توں قضیاتی رابطے $\lnot$، $\land$،
$\lor$، $\lif$ تے $\liff$ ورت کے بندے نیں۔ کلاسیکی منطق دی
معنیات متعین کرن لئی دو صداقتی قدراں $\True$ تے $\False$
ورتدے آں۔ !!{propositional variable}s دی تشریح !!a{valuation}~$\pAssign{v}$
وچ کردے آں۔ ایہ صداقتی قدراں $\True$، $\False$ نوں
!!{propositional variable}s نال جوڑدی اے۔ فیر ہر
!!{valuation} صداقتی قدر $\pValue{v}(!A)$ متعین کردی اے
کسے وی !!{formula}~$!A$ لئی۔ !!^a{formula}،
!!a{valuation}~$\pAssign{v}$ وچ پورا اتردا اے، یعنی $\pSat{v}{!A}$،
اُدوں تے صرف اُدوں جدوں $\pValue{v}(!A) =
\True$۔

کئی قدراں والیاں منطقاں کلاسیکی دو قدری منطق دی وسعت نیں۔
ایہناں وچ $\True$ تے $\False$ توں علاوہ ہور صداقتی قدراں وی
ہو سکدیاں نیں۔ سو کئی قدراں والی منطق وچ !!a{valuation}~$\pAssign{v}$
اوہ فنکشن اے جیہڑا ہر !!{propositional variable}~$p$ نوں ممکن
صداقتی قدراں وچوں اک دیندا اے۔ عام طور اُتے جائز صداقتی
قدراں دے سیٹ نوں~$V$ آکھاں گے۔ کلاسیکی منطق وی کئی قدراں
والی منطق اے جیہدے لئی $V =
\{\True, \False\}$۔ اوہدے وچ صداقتی قدر~$\pValue{v}(!A)$
رابطیاں دے جانے پہچانے صداقتی جدولاں نال کڈھی جاندی اے۔

جدوں ہور صداقتی قدراں شامل کردے آں، تاں اوہناں رابطیاں لئی
~$\pValue{v}(!A)$ کڈھن دے اک توں ودھ سُبھاوک طریقے ہوندے
نیں جنہاں نوں اسیں «تے»، «یا»، «نہیں» تے «جے ... تاں»
پڑھدے آں۔ سو کئی قدراں والی منطق نوں صداقتی قدراں دا سیٹ
تے ہر رابطے لئی ساڈے چُنے ہوئے \emph{صداقتی فنکشن} رل کے
متعین کردے نیں۔ کئی قدراں والی منطق~$\Log L$ لئی ایہ
فنکشن چُنے جان، تاں صداقتی قدر $\pValue{v}(!A)[\Log L]$
قدر بندی نال اکّو طرح متعین ہوندی اے، جیویں کلاسیکی منطق
وچ۔ کلاسیکی منطق وانگ کئی قدراں والیاں منطقاں
\emph{صداقتی فنکشوناں اُتے ٹکیاں} ہوندیاں نیں۔

معنیات دیاں ایہ بنیادی گلاں ہتھ آ جان، تاں اسیں ہمیشہ صادق
ہون، منطقی نتیجہ ہون تے پورا اتر سکّن دے ملدے جُلدے مفہوم
متعین کر سکدے آں۔ کلاسیکی منطق وچ !!a{formula} ہمیشہ صادق
اے جے اوہدی صداقتی قدر $\pValue{v}(!A) = \True$ ہووے ہر~$\pAssign{v}$
لئی۔ کئی قدراں والی منطق وچ اینوں وی کجھ عام کرنا پیندا اے۔
پہلی گل، صداقتی قدراں دے سیٹ~$V$ وچ~$\True$ ہونا ضروری نہیں۔
مثلاً کجھ منطقاں صداقتی قدراں وجوں عدد ورتدیاں نیں، جیویں
$0$ تے~$1$ وچکارلے سارے ناطق عدد۔ اوہ صورت وچ $1$ عام طور
اُتے~$\True$ دا کم کردا اے۔ ہور منطقاں وچ کئی صداقتی قدراں ایہ
کم کردیاں نیں۔ سو ہر کئی قدراں والی منطق لئی اسیں
\emph{مقرر قدراں} دا سیٹ~$V^+$ ضروری منّدے آں۔ فیر آکھ سکدے آں کہ
!!a{formula}، !!a{valuation}~$\pAssign{v}$ وچ پورا اتردا اے،
یعنی $\pSat{v}{!A}[\Log L]$، اُدوں تے صرف اُدوں جدوں
$\pValue{v}(!A)[\Log L] \in V^+$۔ !!^a{formula}~$!A$ منطق
دا ہمیشہ صادق فارمولا اے، یعنی $\Entails[\Log L] !A$،
اُدوں تے صرف اُدوں جدوں $\pValue{v}(!A) \in V^+$ ہووے ہر~$\pAssign{v}$
لئی۔ آخر وچ، آکھدے آں کہ $!A$، !!{formula}s دے سیٹ دا
منطقی نتیجہ اے، یعنی $\Gamma \Entails[\Log L] !A$، جے ہر
!!{valuation} جیہدے وچ~$\Gamma$ والے سارے !!{formula}s پورے
اتردے نیں، اوہدے وچ~$!A$ وی پورا اتردا اے۔

\end{document}
